\centerline{\bf \Huge Гармонические четырёхугольники}

\textbf{Определение.}
Четырёхугольник $A B C D$ называется гармоническим, если он вписанный и для него выполнено $A B \cdot C D=B C \cdot A D$.

\problem{1}
Докажите, что для вписанного четырёхугольника $A B C D$ следующие условия равносильны.

a) $A B C D$ - гармонический. 

б) $B D$ - симедиана треугольника $A B C$.

в) Для точки $M$, середины $A C$, выполнено $\angle A M B=\angle D C B$.

г) Для точки $M$, середины $A C$, выполнено $\angle A M B=\angle A M D$.

д) Касательные в точках $B$ и $D$ к описанной окружности четырёхугольника пересекаются на прямой $A C$ или ей параллельны.

e) Биссектрисы углов $A$ и $C$ пересекаются на диагонали $B D$.

ё) Некое условие на образ $A B C D$ после инверсии относительно одной из вершин - сформулируйте его сами.

\problem{2}
Четырёхугольник $A B C D$ вписан в окружность, причём касательные в точках $B$ и $D$ пересекаются в точке $K$, лежащей на прямой $A C$. Прямая, параллельная $B K$, пересекает прямые $B A, B D, B C$ в точках $P, Q, R$ соответственно. Докажите, что $P Q=Q R$.

\problem{3}
Пусть $P T$ и $P B$ - две касательные к окружности, $A B$ - ее диаметр, и $T H$ перпендикуляр, опущенный из точки $T$ на $A B$. Докажите, что прямая $A P$ делит пополам отрезок $T H$.

\problem{4}
В остроугольном треугольнике $A B C$ на высоте $B K$ как на диаметре построена окружность $\Omega$, пересекающая стороны $A B$ и $B C$ в точках $E$ и $F$ соответственно. К окружности $\Omega$ в точках $E$ и $F$ проведены касательные. Докажите, что их точка пересечения $P$ лежит на прямой, содержащей медиану треугольника $A B C$, проведённую из вершины $B$.

\problem{5}
В окружности $\Omega$ проведены две параллельные хорды $A B$ и $C D$. Прямая, проведённая через $C$ и середину $A B$, вторично пересекает $\Omega$ в точке $E$. Точка $K$ середина отрезка $D E$. Докажите, что $\angle A K E=\angle B K E$.

\newpage

\hspace{3mm}

\problem{6}
Дан параллелограмм $A B C D$. Прямая $l$ перпендикулярна $B C$ и проходит через $B$. Две окружности с общей хордой $C D$ касаются прямой $l$ в точках $P$ и $Q$. Докажите, что отрезки $D P$ и $D Q$ видны из середины $A B$ под равными углами.

\problem{7}
а) В параллелограмме $A B C D$ дана точка $M$ такая, что $\angle M A D=\angle M C D$. Докажите, что $\angle M B A=\angle M D A$.

\noindent
б) В треугольнике $A B C$ дана точка $P$ такая, что $\angle A C P=\angle A B P$. Точка $Q$ симметрична точке $P$ относительно середины $B C$. Докажите, что $\angle C A P=\angle B A Q$.

\noindent
в) Внутри четырёхугольника $A B C D$ дана точка $P$ такая, что $\angle C D P=\angle A B D$, $\angle C B P=\angle A D B$ и $B P=D P$. Докажите, что $P$ лежит на диагонали $A C$.

\problem{8}
Касательные в точках $A$ и $C$ к описанной окружности остроугольного треугольника $A B C$ пересекаются в точке $F$. Внутри треугольника $B F C$ нашлась такая точка $K$, что $\angle K F B=\angle K B C=\angle K C F$. Точка $L$ на стороне $A B$ такова, что $\angle L C B=\angle B F C$. Докажите, что прямая $B K$ делит отрезок $L C$ пополам.